\section{Reinforcement Learning and Post-Training}

\begin{frame}{What Is a Policy?}
  \begin{itemize}\itemsep0.4em
    \item[\textcolor{red}{$\blacktriangleright$}] A \textbf{policy} $\pi$ specifies how an agent (or model) chooses actions.
    \item[\textcolor{red}{$\blacktriangleright$}] \textbf{Deterministic:} $\pi(s) = a$.
    \item[\textcolor{red}{$\blacktriangleright$}] \textbf{Stochastic:} $\pi(a \mid s) = \Pr(a_t = a \mid s_t = s)$.
    \item[\textcolor{red}{$\blacktriangleright$}] For LLMs: $s$ is the prompt/context, $a$ is the next token; the policy is $\pi_\theta(\text{token} \mid \text{context})$.
    \item[\textcolor{red}{$\blacktriangleright$}] \textbf{Goal:} find $\pi^\star$ that maximizes expected cumulative reward (or, in RLHF, human-aligned reward).
  \end{itemize}
\end{frame}

\begin{frame}{Policy Gradients: Learn the Policy Directly}
  \begin{itemize}\itemsep0.35em
    \item[\textcolor{red}{$\blacktriangleright$}] Value-based methods (Q-learning) can be hard in huge state spaces.
    \item[\textcolor{red}{$\blacktriangleright$}] Often easier to parameterize a policy class $\Pi = \{\pi_\theta \mid \theta \in \mathbb{R}^m\}$.
    \item[\textcolor{red}{$\blacktriangleright$}] Maximize expected return $\mathcal{J}(\theta) = \mathbb{E}\big[\sum_{t \geq 0} \gamma^t r_t \mid \pi_\theta\big]$.
    \item[\textcolor{red}{$\blacktriangleright$}] \textbf{Policy gradient:} update $\theta$ along $\nabla_\theta \mathcal{J}(\theta)$ (gradient ascent on parameters).
    \item[\textcolor{red}{$\blacktriangleright$}] PPO and GRPO are policy-gradient-style updates on language-model weights.
  \end{itemize}
\end{frame}

\begin{frame}{Proximal Policy Optimization (Recap)}
    \begin{itemize}\itemsep0.32em
      \item[\textcolor{red}{$\blacktriangleright$}] \textbf{Problem:} large policy updates can collapse performance.
      \item[\textcolor{red}{$\blacktriangleright$}] \textbf{Idea:} take the biggest safe improvement step on the current batch of data.
      \item[\textcolor{red}{$\blacktriangleright$}] \textbf{Clipped objective} (widely used variant): with $r_t(\theta) = \frac{\pi_\theta(a_t \mid s_t)}{\pi_{\theta_k}(a_t \mid s_t)}$,
        \[
          \mathcal{L}^{\mathrm{CLIP}} = \mathbb{E}_t\Big[\min\big(r_t(\theta)\hat{A}_t,\; \mathrm{clip}(r_t(\theta), 1-\epsilon, 1+\epsilon)\hat{A}_t\big)\Big]
        \]
      \item[\textcolor{red}{$\blacktriangleright$}] $\epsilon$ is small (e.g., $0.2$), which limits how far the new policy can drift per update.
    \end{itemize}
  \end{frame}

\begin{frame}{PPO Pros and Cons (Recap)}
  \begin{itemize}\itemsep0.4em
    \item[\textcolor{red}{$\blacktriangleright$}] \textbf{Stable:} clipped updates keep optimization well behaved.
    \item[\textcolor{red}{$\blacktriangleright$}] Encourages exploration of phrasings at each rollout.
    \item[\textcolor{red}{$\blacktriangleright$}] \textbf{Expensive:} many rollouts, critic + RM + reference in memory.
    \item[\textcolor{red}{$\blacktriangleright$}] \textbf{Sensitive:} reward-model mistakes and distribution shift degrade it.
  \end{itemize}
\end{frame}

\begin{frame}{Direct Preference Optimization (Recap)}
  \begin{itemize}\itemsep0.35em
    \item[\textcolor{red}{$\blacktriangleright$}] Bypasses explicit RM training and RL rollouts (Rafailov et al., 2023).
    \item[\textcolor{red}{$\blacktriangleright$}] Optimizes the policy directly from preference pairs $(y^+, y^-)$.
    \item[\textcolor{red}{$\blacktriangleright$}] \textbf{DPO loss} (simplified form):
      \[
        L_{\mathrm{DPO}} = -\log \sigma\Big(\beta \log \tfrac{\pi_\theta(y^+)}{\pi_0(y^+)} - \beta \log \tfrac{\pi_\theta(y^-)}{\pi_0(y^-)}\Big)
      \]
    \item[\textcolor{red}{$\blacktriangleright$}] $\pi_0$ = frozen reference (SFT) policy; $\beta$ = temperature.
    \item[\textcolor{red}{$\blacktriangleright$}] Often simpler and more stable than PPO for alignment, but still preference-based rather than driven by verifiable math or code rewards.
  \end{itemize}
\end{frame}
